%% ============================================================
%% 分片 08：原文第 4 节「Compactification And The Geometry Of Hitchin's Moduli Space」（原文 glshort.tex 起于第 2325 行）
%% 由 tools/harvmac2tex.py 从 arXiv 源码（hep-th/0604151v3）自动生成。
%% 只译英文散文；公式、编号、\label、\cite 一律不改。
%% ============================================================

\section{紧化与 Hitchin 模空间的几何}

\label{sec:compacthitchin}

\begin{keytakeaways}
\begin{itemize}[leftmargin=1.4em,itemsep=2pt,topsep=2pt]
\item \textbf{物理图像}：把四维扭化 ${\EUN}=4$ 超对称杨--米尔斯理论放到 $M=\Sigma\times C$ 上，并令 $C$ 远小于 $\Sigma$；于是在 $\Sigma$ 上只剩一个以 Hitchin 模空间 $\MH(G,C)$ 为靶空间的超对称 sigma 模型，四维规范场论退居为构造 Wilson 算子与 't Hooft 算子的「母体」。
\item \textbf{数学工具}：$C$ 上的 Hitchin 方程 $F-\phi\wedge\phi=0$、$D\phi=D^*\phi=0$；当亏格 $g>1$ 时，$\MH(G,C)$ 是复维数为 $(2g-2){\rm dim}\,G$ 的复流形，它的奇点对应于带有连续未破缺规范对称性的可约解。
\item \textbf{关键结论}：$\MH$ 光滑点处的位形给出从 $\Sigma$ 到 $\MH$ 的常映射（零作用量），近似零作用量则给出缓变映射；$\Sigma$ 远大于 $C$ 的假设保证缓变映射存在且作用量很小，于是 $\Sigma$ 上的长距离物理就由该 sigma 模型描述。
\item \textbf{边界展望}：上述图像在 $\MH$ 的奇点处失效（此时 $\sigma\not=0$，解不必从 $C$ 拉回），$G=U(1)$、$U(N)$ 的情形也需要额外修正；本文因此始终避开奇点，只考虑不可约平坦丛的 Langlands 对应。
\end{itemize}
\end{keytakeaways}

现在我们考虑紧化到两维。我们取四维流形为 $M=\Sigma\times C$，即两个二维流形的乘积。我们取 $C$ 为紧黎曼面，并将在其上研究几何 Langlands 纲领。一般地，我们假设 $C$ 的亏格大于一，从而对单规范群 $G$，$C$ 上的通有平坦联络是不可约的。$\Sigma$ 可以是完备但非紧的二维流形（如 $\Bbb{R}^2$），也可以是第二个紧黎曼面。在后一情形，我们假设 $\Sigma$ 的尺度远大于 $C$ 的尺度。我们的目标是把 $M=\Sigma\times C$ 上的四维规范场论约化为 $\Sigma$ 上一个有效场论的另一种描述。

为求出 $\Sigma$ 上的有效物理，我们必须找出 $M$ 上使作用量（欧几里得号差）或能量（洛伦兹号差）取极小或接近极小的位形。无论用哪种方式，都会得到 \cite{ref18} 中所得的结论（\cite{ref19} 中有密切相关的结果）：四维扭化 ${\EUN}=4$ 超对称规范场论在 $\Sigma$ 上约化为一个 sigma 模型，其映射为 $\Phi:\Sigma\to\MH$，其中 $\MH$ 是 $C$ 上规范群为 $G$ 的 Hitchin 方程解的模空间。

这可由我们在第~\ref{sec:toplag}~节中的结果得到。除拓扑项（它不影响运动方程，也不影响经典位形的能量）之外，扭化 ${\EUN}=4$ 超杨--米尔斯理论的作用量是平方和。由 \eqref{eq:wingo}\ 与 \eqref{eq:uwingo} 可见，它恰在所有 $t$ 下拓扑方程都被满足时取极小，换言之，当 ${\cal F}=D^*\phi=0$，同时 $D\sigma=[\phi,\sigma]=[\sigma,\bar\sigma]=0$ 时取极小。这些方程的一类简单解是取 $A$ 与 $\phi$ 为来自 $C$ 的拉回，且 $\sigma=0$。事实上，对黎曼面 $C$ 上的联络 $A$ 与伴随值 1-形式 $\phi$ 而言，方程 ${\cal F}=D^*\phi=0$ 就是 Hitchin 方程 \cite{ref20} 的一种书写方式。取 ${\cal F}$ 的实部与虚部，这些方程也可改写为
% 式 (4.1)
\setcounter{equation}{0}
\begin{equation}\label{eq:zongo}
\begin{aligned} F-\phi\wedge\phi & = 0\\ D\phi=D^*\phi&=0
.\\\end{aligned}
\end{equation}

我们将这些方程在 $C$ 上的解模去规范变换所得的模空间记作 $\MH$，或更详细地记作 $\MH(G,C)$。若 $C$ 的亏格 $g$ 大于一，则 $\MH(G,C)$ 是复维数为 $(2g-2){\rm dim}\,G$ 的复流形。它带有奇点，对应于 Hitchin 方程的可约解，即具有连续未破缺规范对称性的解。

正如我们在第~\ref{sec:vanishing}~节末尾所解释的，$M$ 上方程 ${\cal F}=D^*\phi=0$ 的任一解都联系着一个从 $\Sigma$ 到 $\MH$ 的映射。此外，若避开 $\MH$ 的奇点，则 $M$ 上的解实际上是从 $C$ 上的解拉回而来。此时，关于 $\sigma$ 的条件要求 $\sigma=0$。因此，一个能量或作用量为零且避开 $\MH$ 奇点的场，就是从 $\Sigma$ 到 $\M_H$ 中某个光滑点的常映射。

相应地，能量或作用量近乎为零的场来自从 $\Sigma$ 到 $\M_H$ 的缓变映射。（若一个映射只在 $\Sigma$ 中远大于 $C$ 的尺度上才有显著变化，就称它是缓变的。）于是，只要不遇到 $C$ 上具有连续规范对称性的场，$\Sigma$ 上长距离处的有效物理就由一个超对称 sigma 模型给出，其靶空间为 $\M_H(G,C)$。

假设 $\Sigma$ 远大于 $C$ 的重要性在于：它保证非常值但缓变的映射 $\Sigma\to \MH(G,C)$ 可以存在，而且还保证决定缓变映射 $\Phi:\Sigma\to \MH(G,C)$ 的一对场 $(A,\phi)$ 具有很小的作用量。\footnote{反之亦然，我们将在第~\ref{sec:wilbranes}~节中说明。给定缓变映射 $\Phi:\Sigma\to \MH(G,C)$，可以重构出 $M$ 上决定该映射且作用量很小的一对场 $(A,\phi)$。} 反之，若假设 $C$ 远大于 $\Sigma$，则代表非常值通有映射 $\Phi:\Sigma\to \MH(G,C)$ 的一对 $(A,\phi)$ 具有很大的作用量；此时，我们反而可以得到一个涉及缓变映射 $\tilde\Phi:C\to \MH(G,\Sigma)$ 的良好描述。

以上分析假设 $G$ 是半单的。若 $G=U(1)$ 或 $U(N)$，则需要作一些修正，因为 Hitchin 方程的每个解都带有一个 $U(1)$ 对称群。因此，低能描述涉及一个 sigma 模型与一个超对称 $U(1)$ 规范场论的乘积。此外，对任意 $G$，低能描述在 $\M_H$ 的奇点处失效，这些奇点对应于具有连续规范对称性（对 $G=U(1)$ 或 $U(N)$ 则是额外的连续规范对称性）的 Hitchin 方程解。在这些点处，额外的自由度开始起作用：能量或作用量为零的场可以有 $\sigma\not=0$，并且不必是从 $C$ 拉回而来。（它具有更一般的分解 $E=\oplus_{i=1}^nE'_i\otimes E''_i$，这一点已在第~\ref{sec:vanishing}~节末尾描述过。）对半单 $G$，若 $C$ 的亏格至少为 2，Hitchin 方程的通有解是不可约的。在本文中，我们将远离 $\M_H$ 的奇点（例如通过考虑不可约平坦 $^L\negthinspace G_{\Bbb{C}}$ 丛的 Langlands 对应）。一些简单的推广将在 \cite{ref26} 中考虑。

以 $\MH(G,C)$ 为靶空间的 sigma 模型将被证明是理解几何 Langlands 纲领的有力工具。但它与底层规范场论的关系很重要。一方面，故事中的关键角色是 Wilson 算子与 't Hooft 线算子，它们最自然地在底层的四维规范场论中构造。此外，当 sigma 模型由于 $\MH$ 的奇点而失效时，底层规范场论是一个重要的工具。但许多东西可以直接在 sigma 模型中理解。

本节余下部分将回顾 $\MH$ 的经典几何，只聚焦于本文所需的那些主题。

\begin{yizhu}[物理图像]
这里的「紧化到两维」是 Kaluza--Klein 式的红外约化：$C$ 的尺度远小于 $\Sigma$，于是 $C$ 上高阶 Kaluza--Klein 模式的能标约为 $1/{\rm size}(C)$，在 $\Sigma$ 的长距离下被积掉，只剩沿 $C$ 不变的零模与缓变模。Hitchin 模空间之所以出现，正因为它参数化了四维拓扑扭化理论在 $C$ 上的 BPS 位形。
\end{yizhu}

\begin{yizhu}[易错点]
式 (4.1) 中 $D^*\phi$ 的星号是 $C$ 上二维的 Hodge 星算子；$D\phi=0$ 与 $D^*\phi=0$ 是两个独立的实方程，不要把它误认成四维的 $D\star\phi=0$。原式第二行末尾的句点是排版遗留，与数学内容无关。
\end{yizhu}

\begin{yizhu}[含糊点拨]
复维数 $(2g-2)\dim G$ 可以拆开来看：稳定 $G$-丛模空间的复维数为 $(g-1)\dim G$，而 Higgs 场 $\phi\in H^0(C,K_C\otimes{\rm ad}E)$ 由 Riemann--Roch 也给出 $(g-1)\dim G$ 个复自由度（$\deg K_C=2g-2$，且稳定丛的 $H^1$ 消失）；两者相加即得。$\MH$ 作为实流形的维数是这个复维数的两倍。
\end{yizhu}

\begin{yizhu}[物理图像]
奇点处 $\sigma\not=0$ 并非细节：它意味着「零作用量位形必为从 $C$ 的拉回」这一论证失效，额外的连续规范对称性带来了新的自由度。这正是本文后面必须回到四维规范场论去处理更一般边界条件的原因，也是后来把 $\MH$ 换成含 $\sigma$ 的 DG 栈的出发点。
\end{yizhu}
